\section{Background}

The global livestock sector is important for supporting food production worldwide, and cattle are the most economically valuable livestock species in this industry. According to the Food and Agriculture Organization (FAO), the global cattle population is over one billion, which is making a significant contribution in producing meat and dairy products \cite{fao2018livestock}. Because of this, farm sustainability, profitability, and productivity rely directly on the health and welfare of these animals. Early detection of health problems like poor body condition, lameness, and abnormal behavior can prevent significant economic losses and improve welfare standards \cite{weary2009understanding}.

One of the most reliable methods that have been practiced for years by farmers in checking the condition of their cattle is the Body Condition Score (BCS). In the past, BCS was calculated by the farmer by visually examining and physically checking the fat deposition in different body parts of the cattle through a score of 1 to 5, increasing in steps of 0.25 units \cite{edmonson1989bcs}. BCS will ensure effective milk production, high fertility rate and resistance to disease infection. However, the problem with the conventional method is that it consumes much more time and labor, especially in a large scale farm. One of the major problems among dairy cows is lameness. Up to 25\% of dairy cows in intensive farming suffer from lameness \cite{whay2003lameness}. The disease causes pain, decreased feed consumption, decreased milk production, and impairs reproduction, but early detection of an abnormal gait pattern is helpful in providing farmers with treatment at an earlier stage.

Automated behavioral analysis and cow identification play a significant role in detecting the health, stress levels, and wellbeing of the animal. Changes in behavior, in relation to feeding, ruminating, lying and social behavior, are recognized before any clinical symptoms of illness become visible \cite{weary2009understanding}. The computer vision technology allows monitoring animals in a non-invasive way and represents a highly reliable method that decreases the necessity of direct observations. Ear tagging and RFID-based identification techniques require manual handling and can be easily misplaced or destroyed. On the other hand, visual recognition systems are not only painless but also easy to incorporate in automated animal health tracking systems.

Recent advancements in computer vision, deep learning, and IoT have opened up great chances for automated cattle health monitoring \cite{neethirajan2020digitalization}. CNN models have proven themselves effective at image classification and object detection, whereas transfer learning allows for using pre-trained networks on large datasets such as ImageNet for agricultural tasks where data is limited. Nevertheless, current solutions in cattle monitoring are mostly limited to executing single tasks. A literature review comprising 33 primary scientific sources performed in the course of this thesis reveals that multi-task cattle monitoring solutions are relatively rare. This negatively impacts the efficiency of automation and effectiveness of the related tasks by depriving them of the benefit of sharing learned features.


\section{Rationale of the Study and Motivation}

There are some problems in existing cattle health monitoring frameworks, which motivate the current research, and also there is an increasing need for precision livestock farming technologies. It is seen from the literature review that around 88\% (29 out of 33) have used private datasets, and therefore, replication cannot be achieved easily. Apart from this, previous studies have only been able to concentrate on one task at a time such as BCS, lameness detection, or cattle identification. The proposed method is not utilizing Multi-Task Learning (MTL), which allows multiple correlated tasks to share representations and thereby save computation costs while improving performance. This is because the two tasks are naturally related in that the physical well-being, such as lameness, influences feeding behavior and affects body condition.

Several published systems depend on costly devices like thermal cameras, wearable sensors, and depth cameras, making large-scale deployment difficult. This highlights the importance of affordable monitoring systems that rely on RGB cameras and can also use depth information if needed. So, Farms often need on-site processing since cloud systems and weak internet connections can slow down monitoring. Although edge computing is effective, many existing architectures are too computationally expensive for edge hardware. Only a small number of cattle monitoring systems use self-supervised learning. Most commercial farms record huge amounts of daily video data using 24/7 cameras. Even though this video data is mostly unused, it can support model pre-training and improve learning with limited labeled data. Overall, these issues demonstrate the need for a lightweight and integrated framework for cattle monitoring.


\section{Problem Statement}

Automated cattle health monitoring has been around for years, yet fundamental problems still remain difficult to resolve. Because of this, most computer vision and deep learning solutions have yet to achieve real commercial deployment, or remain limited in practical use when they do.

The existing systems handle each of the monitoring tasks separately, which requires multiple models to cover a full health assessment. Running multiple separate models increases computational load, prevents the related tasks from sharing learned features, complicates the system deployment and ignores the connections between health indicators. Looking at existing literature shows that there is no framework that is capable of addressing body condition scoring, behavior recognition, lameness detection, and individual identification together. While multi-task learning is often mentioned in research, few studies have actually built and tested fully integrated systems.

In order to properly validate the results, compare different methods fairly and to build on existing works, standardized dataset evaluation is essential. However most cattle monitoring studies rely solely on datasets which are proprietary and not publicly available. The absence of openly shared datasets makes it very difficult to reproduce results, compare the model performance, and limits how much the field can grow over time. This is very different from general computer vision which has benefited greatly from COCO or ImageNet.

Beyond this, most studies evaluate their models using data from a single farm or cattle breed, which makes it hard to know how well these methods would work in different conditions. Testing under such specific conditions can easily make models overfit to particular farm settings, camera angles, lighting and may not transfer across different breeds. Seasonal change, changes across environment and time is also largely ignored while evaluating these models. Most importantly, cross-dataset evaluation, where a model trained on one dataset is tested on a completely different one, is rarely done in existing studies.

Lastly, depth and thermal imaging have both gained considerable interest for cattle monitoring, however lack of open datasets remains a major limitation to multimodal research in this area of study. No public thermal datasets are available, only MmCows provides IMU data, and depth data is offered by both MmCows and Dryad. This shortage of public datasets leaves researchers with two options --- collect their own private dataset or give up on multimodal approaches.


\section{Objectives}

The primary objective of this research is to develop a lightweight, edge-deployable, and openly evaluated multi-task deep learning framework for automated cattle health and behavior monitoring. By leveraging public datasets and shared feature representations, the framework aims to concurrently perform multiple health-related assessments while establishing a standardized baseline.

To achieve this goal, the research sets out the following specific objectives:

\begin{enumerate}
    \item \textbf{To design a unified, lightweight multi-task architecture} with a shared backbone and specialized task-specific heads, and to evaluate sequential training strategies to optimize performance across Body Condition Scoring, behavior recognition, lameness detection, and individual identification.
    
    \item \textbf{To improve reproducibility} by conducting all evaluations exclusively on publicly available datasets and documenting preprocessing, task-specific splits, model settings, and evaluation protocols.
    
    \item \textbf{To conduct preliminary cross-dataset evaluation} to identify domain shift challenges across different farms and environmental conditions, establishing a baseline for comprehensive generalization studies in future work.
\end{enumerate}

These objectives address the reproducibility gap in agricultural AI research and create a scalable foundation for real-time cattle health monitoring across various farm environments.


\section{Methodology in Brief}

The proposed methodology follows a structured, multi-phase process to design, implement, and validate the multi-task deep learning framework. This research workflow is divided into five key phases: data preparation, architecture design, sequential task training, evaluation, and ablation studies.


\subsection{Evaluation Protocol}

Evaluation metrics are carefully tailored to each task to address their distinct outputs and operational risks. For body condition scoring, the model's performance is measured using Mean Absolute Error (MAE) alongside classification accuracy within $\pm0.25$ and $\pm0.5$ tolerances. For behavior recognition, evaluation metrics include per-class accuracy, macro-averaged F1-score, and confusion matrix analysis. For lameness detection, evaluation metrics include accuracy, F1-score, and AUC-ROC. Individual identification is assessed using Top-1 accuracy. Finally, cross-dataset evaluation is conducted by training on a primary dataset and testing on an independent dataset to rigorously assess the model's generalization capabilities.


\subsection{Ablation Studies}

A series of structured ablation experiments will be conducted to evaluate the contributions of key framework components. For backbone selection, EfficientNet-B0 will be compared against MobileNetV3-Small and ResNet-50 to determine the optimal trade-off between computational efficiency and accuracy. For the body condition scoring task, the performance of the Rank Consistent Ordinal Regression (CORAL) loss will be compared against standard Cross-Entropy loss. Regarding training strategies, experiments will compare sequential learning and joint multi-task training. The impact of multi-modal data will be discussed through a cross-dataset modality observation. Finally, the efficacy of attention mechanisms will be evaluated by comparing architectures with and without the Convolutional Block Attention Module (CBAM).


\section{Scope and Challenges}

\subsection{Scope of the Research}

To ensure feasibility and depth of investigation, this research is defined by the following boundaries:

\begin{enumerate}
    \item \textbf{Four Core Tasks:} The scope is restricted to Body Condition Scoring, behavior recognition, lameness detection, and individual identification, selected due to their practical significance and public data availability.
    
    \item \textbf{Public Datasets Only:} All experimental evaluations are conducted exclusively on publicly available datasets to ensure a more transparent and reproducibility-oriented framework, using public datasets and documented preprocessing/splits.
    
    \item \textbf{RGB as Primary Modality:} RGB imagery serves as the primary input modality, with depth information evaluated solely as an auxiliary modality for the body condition scoring task using the Dryad dataset.
    
    \item \textbf{Dairy Cattle Focus:} The study focuses primarily on Holstein dairy cattle, which represent the dominant breed in public cattle monitoring datasets.
    
    \item \textbf{Lightweight Architectures:} To ensure suitability for resource-constrained edge deployment, the backbone selection is restricted to architectures with fewer than 10 million parameters.
    
    \item \textbf{Sequential Task Training:} Frozen-backbone sequential training serves as the primary optimization strategy, while joint multi-task fine-tuning and loss balancing are explored as secondary extensions.
\end{enumerate}


\subsection{Challenges}

\textbf{Camera Geometry Conflicts}

The four tasks in this framework require fundamentally different camera viewpoints: BCS scoring performs best from a rear-view angle (capturing hip and spine morphology), while Lameness detection requires a side-view angle (to observe gait stride length and limb symmetry). Running both tasks from a single fixed camera creates an unavoidable geometric trade-off. To address this, the deployment model assumes a multi-camera setup where a single edge device loads the shared backbone once, but selectively activates the rear-view or side-view task heads based on the incoming camera feed.



\textbf{Dataset Heterogeneity}

Public cattle datasets exhibit high heterogeneity, including variations in image resolution, quality, camera placement (top-view, rear-view, and side-view), annotation protocols, breed characteristics, and farm environments. To mitigate these discrepancies, the methodology incorporates a standardized preprocessing pipeline, robust data augmentation, and domain adaptation techniques.


\textbf{Class Imbalance}

Cattle monitoring datasets frequently exhibit severe class imbalances. In lameness detection, lame individuals are heavily outnumbered by healthy cattle; similarly, extreme body condition scores (under 2.5 or over 3.5) are significantly rarer than moderate scores, and specific behaviors (such as drinking or licking) occur infrequently. To address these imbalances, the framework employs Focal Loss for classification tasks, class-weighted sampling during training, and balanced evaluation metrics.


\textbf{Temporal Modeling for Video Tasks}

While behavior recognition and lameness detection benefit from temporal context, video-based modeling increases computational overhead and introduces architectural complexity (e.g., through LSTMs or I3D networks), alongside potential discrepancies between frame-level and clip-level annotations. The framework addresses this by first utilizing frame-level classification combined with temporal aggregation, treating recurrent model integration as a subsequent extension.


\textbf{Preventing Negative Transfer}

Multi-task learning designs are susceptible to negative transfer, where optimization for one task degrades the performance of another. This typically occurs due to conflicting task gradients during joint optimization, varying optimal feature representations across tasks, and disparity in task difficulties. The primary mitigation strategy involves sequential training with a frozen backbone to decouple task optimization, supplemented by gradient normalization and dynamic loss weighting during joint fine-tuning.


\textbf{Limited Labeled Data}

Despite utilizing multiple public datasets, the total volume of labeled data remains small compared to standard computer vision benchmarks. For example, the ScienceDB BCS dataset contains 53,566 images, in contrast to the millions in ImageNet. The MmCows dataset tracks only 16 individuals over a two-week period, and the CattleLameness dataset contains a relatively small number of lame examples. To mitigate data limitations, the methodology leverages transfer learning from ImageNet, robust regularization (such as dropout and weight decay), and extensive data augmentation, with self-supervised pre-training on unlabeled video frames designated as an optional extension.
